% from ega0-1-fr.tex
\oldpage[0\textsubscript{I}]{32}

\begin{env}[3.5.4]
\label{0.3.5.4-fr}
Supposons maintenant que la catégorie $\mathbf{K}$ soit telle\footnote{Dans
le livre cité dans l'Introduction, nous donnerons des conditions très
générales sur la catégorie $\mathbf{K}$ assurant l'existence des images
réciproques de préfaisceaux à valeurs dans $\mathbf{K}$.} que \emph{tout}
préfaisceau $\mathcal{G}$ sur $Y$ à valeurs dans $\mathbf{K}$ admette une
image réciproque par $\psi$, que nous noterons $\psi^*(\mathcal{G})$.

Nous allons voir qu'on peut définir $\psi^*(\mathcal{G})$ comme \emph{foncteur
covariant} en $\mathcal{G}$, de la catégorie des préfaisceaux sur $Y$ à
valeurs dans $\mathbf{K}$, dans celle des faisceaux sur $X$ à valeurs dans
$\mathbf{K}$, de telle façon que l'isomorphisme $v\to v^b$ soit un
\emph{isomorphisme de bifoncteurs}
\begin{equation*}
  \operatorname{Hom}_X(\psi^*(\mathcal{G}),\mathcal{F})
  \xrightarrow{\sim}
  \operatorname{Hom}_Y(\mathcal{G},\psi_*(\mathcal{F}))
\tag{3.5.4.1}\label{0.3.5.4.1-fr}
\end{equation*}
en $\mathcal{G}$ et $\mathcal{F}$.

En effet, pour tout morphisme
$w:\mathcal{G}_1\to\mathcal{G}_2$ de préfaisceaux sur $Y$ à valeurs dans
$\mathbf{K}$, considérons le morphisme composé
\[
  \mathcal{G}_1\xrightarrow{w}\mathcal{G}_2
  \xrightarrow{\rho_{\mathcal{G}_2}}
  \psi_*(\psi^*(\mathcal{G}_2))~;
\]
il lui correspond un morphisme
$(\rho_{\mathcal{G}_2}\circ w)^\sharp:
\psi^*(\mathcal{G}_1)\to\psi^*(\mathcal{G}_2)$, que nous noterons
$\psi^*(w)$. On a donc, en vertu de (3.5.3.3)
\begin{equation*}
  \psi_*(\psi^*(w))\circ\rho_{\mathcal{G}_1}
  =\rho_{\mathcal{G}_2}\circ w.
\tag{3.5.4.2}\label{0.3.5.4.2-fr}
\end{equation*}
Pour tout morphisme $u:\mathcal{G}_2\to\psi_*(\mathcal{F})$, où
$\mathcal{F}$ est un faisceau sur $X$ à valeurs dans $\mathbf{K}$, on a,
d'après (3.5.3.3), (3.5.4.2) et la définition de $v^b$
\[
  (u^\sharp\circ\psi^*(w))^b
  =\psi_*(u^\sharp)\circ\psi_*(\psi^*(w))\circ\rho_{\mathcal{G}_1}
  =\psi_*(u^\sharp)\circ\rho_{\mathcal{G}_2}\circ w
  =u\circ w
\]
ou encore
\begin{equation*}
  (u\circ w)^\sharp=u^\sharp\circ\psi^*(w).
\tag{3.5.4.3}\label{0.3.5.4.3-fr}
\end{equation*}
Si on prend en particulier pour $u$ un morphisme
\[
  \mathcal{G}_2\xrightarrow{w'}\mathcal{G}_3
  \xrightarrow{\rho_{\mathcal{G}_3}}
  \psi_*(\psi^*(\mathcal{G}_3)),
\]
il vient
\[
  \psi^*(w'\circ w)
  =(\rho_{\mathcal{G}_3}\circ w'\circ w)^\sharp
  =(\rho_{\mathcal{G}_3}\circ w')^\sharp\circ\psi^*(w)
  =\psi^*(w')\circ\psi^*(w),
\]
d'où notre assertion.

Enfin, pour tout faisceau $\mathcal{F}$ sur $X$ à valeurs dans
$\mathbf{K}$, soit $i_{\mathcal{F}}$ le morphisme identique de
$\psi_*(\mathcal{F})$ et notons
\[
  \sigma_{\mathcal{F}}:
  \psi^*(\psi_*(\mathcal{F}))\to\mathcal{F}
\]
le morphisme $(i_{\mathcal{F}})^\sharp$~; la formule (3.5.4.3) donne en
particulier la factorisation
\begin{equation*}
  u^\sharp:\psi^*(\mathcal{G})
  \xrightarrow{\psi^*(u)}
  \psi^*(\psi_*(\mathcal{F}))
  \xrightarrow{\sigma_{\mathcal{F}}}\mathcal{F}
\tag{3.5.4.4}\label{0.3.5.4.4-fr}
\end{equation*}
pour tout morphisme $u:\mathcal{G}\to\psi_*(\mathcal{F})$. Nous dirons que le
morphisme $\sigma_{\mathcal{F}}$ est \emph{canonique}.
\end{env}

\begin{env}[3.5.5]
\label{0.3.5.5-fr}
Soit $\psi':Y\to Z$ une application continue, et supposons que tout
préfaisceau $\mathcal{H}$ sur $Z$ à valeurs dans $\mathbf{K}$ admette une
image réciproque $\psi'^*(\mathcal{H})$ par $\psi'$. Alors (avec les
hypothèses de (3.5.4)) tout préfaisceau $\mathcal{H}$ sur $Z$ à valeurs dans
$\mathbf{K}$ admet une image réciproque par $\psi''=\psi'\circ\psi$ et l'on a
un isomorphisme canonique fonctoriel
\begin{equation*}
  \psi''^*(\mathcal{H})\xrightarrow{\sim}
  \psi^*(\psi'^*(\mathcal{H}))
\tag{3.5.5.1}\label{0.3.5.5.1-fr}
\end{equation*}

